\section{常用的线性结构}

\begin{frame}
    \frametitle{\insertsectionhead：引例一}
    \begin{exampleblock}{编辑器撤销}
        用户依次输入$A$、$B$、$C$，随后连续执行两次撤销。
        应先撤销哪次操作？如何保存尚未撤销的操作？
    \end{exampleblock}
    \begin{block}{关键需求}
        撤销总是处理最近的一次操作，需要后进先出的记录方式。
    \end{block}
\end{frame}

\subsection{栈}

\begin{fragile}
    \frametitle{\insertsubsectionhead}
    \bicolumns[0.6]{
        \begin{block}{栈(Stack)}
            \begin{itemize}
                \item 一种操作\alert{受限}的序列
                \item \alert{仅}允许在一端插入删除元素
                \item 可操作一端为\textbf{栈顶(top)}，另一端为\textbf{栈底(bottom)}
                \item \textbf{后进先出(Last In First Out, LIFO)}
            \end{itemize}
        \end{block}
        \uncover<2>{
            \begin{exampleblock}{日常生活中的栈}
                \begin{itemize}
                    \item 一摞碗、碟、凳子等
                    \item 糖葫芦、牛羊肉串等
                    \item 子弹夹
                    \item $\dots$
                \end{itemize}
            \end{exampleblock}
        }
    }{
        \begin{figure}
            \centering
            \begin{bytefield}{6}
                \stackbox{4.5}{6}{$\vdots$\\\vspace*{1ex}$\uparrow\downarrow$} \\
                \begin{leftwordgroup}{$n$}
                    \begin{rightwordgroup}{\mintinline[fontsize=\smaller]{c}{栈顶}}
                        \stackbox{1.5}{6}{$a_{n}$}
                    \end{rightwordgroup} \\
                    \stackbox{2.5}{6}{$\vdots$\vspace{1.5ex}} \\
                    \stackbox{1.5}{6}{$a_{2}$} \\
                    \begin{rightwordgroup}{\mintinline[fontsize=\smaller]{c}{栈底}}
                        \stackbox{1.5}{6}{$a_{1}$}
                    \end{rightwordgroup}
                \end{leftwordgroup}
            \end{bytefield}
            \caption{栈的示意图}
            \label{fig:demo_stack}
        \end{figure}
    }
\end{fragile}

\begin{fragile}
    \frametitle{\insertsectionhead}
    \begin{block}{栈的抽象数据类型}
        \begin{minted}[linenos=false,escapeinside=@@]{c}
            ADT Stack {
            数据:
                数据对象: @$\mathcal{D} = \{a_{k} | a_{k}\in\text{数据元素集合}, k\in\mathbb{Z}\cap[1,n]\}$@
                逻辑关系: @$\mathcal{R} = \{\langle{}a_{k-1},a_{k}\rangle | k\in\mathbb{Z}\cap[2,n]\}$@
            操作:
                create_stack()
                    构造并初始化一个空栈@$s$@
                is_empty_stack()
                    判断栈是否为空
                push_stack(s, x)
                    若栈@$s$@存在且未满, 则在其栈顶插入值为@$x$@的新元素
                pop_stack(s)
                    若栈@$s$@存在且非空, 则删除其栈顶元素
                top_stack(s)
                    若栈@$s$@存在且非空, 则返回其栈顶元素
            }
        \end{minted}
    \end{block}
\end{fragile}

\begin{fragile}
    \frametitle{\insertsubsectionhead}
    \bicolumns[0.5]{
        \begin{block}{栈的类型说明}
            \begin{itemize}
                \item 可利用已有结构的再封装
                \item \alert{顺序栈}可采用序列
                \item \alert{链式栈}可采用链表
            \end{itemize}
        \end{block}
    }{
        \begin{minted}[highlightlines={}]{c}
            typedef struct {
                SequenceList *_;
            } SequenceStack;
        \end{minted}
        \begin{minted}[highlightlines={}]{c}
            typedef struct {
                LinkedList *_;
            } LinkedStack;
        \end{minted}
    }
\end{fragile}

\begin{fragile}
    \frametitle{\insertsubsectionhead}
    \bicolumns[0.5]{
        \begin{block}{栈的初始化}
            \begin{itemize}
                \item 可利用序列/链表已有的初始化方法
                \item 注意空指针的处理
            \end{itemize}
        \end{block}
    }{
        \begin{minted}[highlightlines={}]{c}
            SequenceStack *create_sequence_stack() {
                SequenceStack *s = malloc(sizeof(SequenceStack));
                if (s) {
                    s->_ = create_sequence_list();
                }
                return s;
            }
        \end{minted}
        \begin{minted}[highlightlines={}]{c}
            LinkedStack *create_linked_stack() {
                LinkedStack *s = malloc(sizeof(LinkedStack));
                if (s) {
                    s->_ = create_linked_list();
                }
                return s;
            }
        \end{minted}
    }
\end{fragile}

\begin{fragile}
    \frametitle{\insertsubsectionhead}
    \bicolumns[0.5]{
        \begin{block}{判断栈是否为空}
            \begin{itemize}
                \item 顺序栈可利用序列已有的判空方法
                \item 链式栈需检测首结点的下一个结点
                \item 此处利用了短路求值原则，下同
            \end{itemize}
        \end{block}
    }{
        \begin{minted}[highlightlines={}]{c}
            bool empty_sequence_stack(SequenceStack *s) {
                return !s || empty_sequence_list(s->_);
            }
        \end{minted}
        \begin{minted}[highlightlines={}]{c}
            bool empty_linked_stack(LinkedStack *s) {
                return !s || !s->_->head->next;
            }
        \end{minted}
    }
\end{fragile}

\begin{fragile}
    \frametitle{\insertsubsectionhead}
    \bicolumns[0.5]{
        \begin{block}{入栈}
            \begin{itemize}
                \item 将新元素\alert{压入}栈中
                \item 指将新元素放入栈顶并更新栈顶
                \item 顺序栈栈顶序号为\mintinline[fontsize=\smaller]{c}{last}
                \item 链式栈栈顶为首结点的直接后继
            \end{itemize}
        \end{block}
    }{
        \begin{minted}[highlightlines={}]{c}
            bool push_sequence_stack(
                    SequenceStack *s, DataType d) {
                return s && insert_sequence_list(
                        s->_, s->_->last + 1, d);
            }
        \end{minted}
        \begin{minted}[highlightlines={}]{c}
            bool push_linked_stack(
                    LinkedStack *s, DataType d) {
                return s && insert_after_linked_by_index(
                        s->_, 0, d);
            }
        \end{minted}
    }
\end{fragile}

\begin{fragile}
    \frametitle{\insertsubsectionhead}
    \bicolumns[0.5]{
        \begin{block}{出栈}
            \begin{itemize}
                \item 将栈顶元素从栈中\alert{弹出}
                \item 指将栈顶元素取出并更新栈顶
            \end{itemize}
        \end{block}
    }{
        \begin{minted}[highlightlines={}]{c}
            bool pop_sequence_stack(
                    SequenceStack *s, DataType *p) {
                return s && remove_sequence_list(
                        s->_, s->_->last, p);
            }
        \end{minted}
        \begin{minted}[highlightlines={}]{c}
            bool pop_linked_stack(
                    LinkedStack *s, DataType *p) {
                return s && remove_linked_by_index(
                        s->_, 1, p);
            }
        \end{minted}
    }
\end{fragile}

\begin{fragile}
    \frametitle{\insertsubsectionhead}
    \bicolumns[0.5]{
        \begin{block}{取栈顶元素}
            \begin{itemize}
                \item 将栈顶元素\alert{复制}一份
                \item 不改变栈
                \item 首先判断栈是否为空
            \end{itemize}
        \end{block}
    }{
        \begin{minted}[highlightlines={}]{c}
            bool top_sequence_stack(
                    SequenceStack *s, DataType *p) {
                if (empty_sequence_stack(s) || !p) {
                    return false;
                }
                *p = s->_->data[s->_->last];
                return true;
            }
        \end{minted}
        \begin{minted}[highlightlines={}]{c}
            bool top_linked_stack(
                    LinkedStack *s, DataType *p) {
                if (empty_linked_stack(s) || !p) {
                    return false;
                }
                *p = s->_->head->next->data;
                return true;
            }
        \end{minted}
    }
\end{fragile}

\begin{fragile}
    \frametitle{\insertsubsectionhead}
    \begin{exampleblock}{栈与递归}
        \begin{itemize}
            \item \textbf{递归(recursion)}：函数直接或间接调用自身的过程
            \item 递归\textbf{要素}：终止条件与递归体
            \item 递归\textbf{实现}：在\alert{栈}中记录每次调用的有用信息
        \end{itemize}
    \end{exampleblock}
    \pause
    \begin{exampleblock}{例：求阶乘}
        \vspace{2ex}
        \bicolumns[0.5]{
            \vspace*{1.5ex}
            \[
                n!=\begin{cases}
                    1            & n=0 \\
                    n\cdot(n-1)! & n>0
                \end{cases}
            \]
        }{
            \begin{minted}{c}
                int factorial(int n) {
                    return 0 == n ? 1 : n * factorial(n - 1);
                }
            \end{minted}
        }
    \end{exampleblock}
\end{fragile}

\begin{fragile}
    \frametitle{\insertsubsectionhead}
    \bicolumns[0.5]{
        \begin{figure}
            \centering
            \begin{bytefield}{2}
                \stackbox{3}{14}{\only<5>{
                        \mintinline[fontsize=\smaller]{console}{factorial()} \\
                        \mintinline[fontsize=\scriptsize]{console}{return address = 1002} \\
                        \mintinline[fontsize=\smaller]{console}{n = 0, ...}
                    }
                } \\
                \stackbox{3}{14}{\only<4-6>{
                        \mintinline[fontsize=\smaller]{console}{factorial()} \\
                        \mintinline[fontsize=\scriptsize]{console}{return address = 1002} \\
                        \mintinline[fontsize=\smaller]{console}{n = 1, ...}
                    }
                } \\
                \stackbox{3}{14}{\only<3-7>{
                        \mintinline[fontsize=\smaller]{console}{factorial()} \\
                        \mintinline[fontsize=\scriptsize]{console}{return address = 2003} \\
                        \mintinline[fontsize=\smaller]{console}{n = 2, ...}
                    }
                } \\
                \stackbox{2}{14}{\only<2-8>{
                        \mintinline[fontsize=\smaller]{console}{main()}\\
                        \mintinline[fontsize=\smaller]{console}{n = 2, ...}
                    }
                }
            \end{bytefield}
            \caption{函数调用栈：\only<2-5>{递进调用}\only<6-9>{回归求值}}
            \label{fig:function_call_stack}
        \end{figure}
    }{
        \begin{minted}[firstnumber=1001,highlightlines={1002}]{c}
            int factorial(int n) {
                return 0 == n ? 1 : n * factorial(n - 1);
            }
        \end{minted}
        \begin{minted}[firstnumber=2001,highlightlines={2003}]{c}
            int main(void) {
                int n = 2;
                int f = factorial(n);
                return 0;
            }
        \end{minted}
    }
\end{fragile}

\begin{frame}
    \frametitle{\insertsubsectionhead}
    \begin{exampleblock}{递归应用：汉诺塔(Hanoi)问题}
        \begin{itemize}
            \item 相传世界之初有钻石宝塔甲，其上有$64$金碟，由大到小依次向上摆放
            \item 附近另有二空塔：乙、丙，与甲类似
            \item 婆罗门牧师试图将甲塔金碟移至丙塔，但需借乙塔中转
            \item 每次只移最上一碟，且不可将大碟置于小碟之上
            \item 牧师完成之日便是世界末日
        \end{itemize}
    \end{exampleblock}
\end{frame}

\begin{fragile}
    \frametitle{\insertsubsectionhead}
    \bicolumns[0.5]{
        \begin{exampleblock}{递归解题思路}
            \begin{enumerate}
                \item 若剩余碟数$n=0$，则结束
                \item 否则执行：
                      \begin{itemize}
                          \item 将$n-1$碟从甲借丙移至乙
                          \item 将甲中剩余一碟从甲移至丙
                          \item 将$n-1$碟从乙借甲移至丙
                      \end{itemize}
            \end{enumerate}
        \end{exampleblock}
    }{
        \begin{minted}{c}
            void move(char a, char b) {
                printf("%c -> %c\n", a, b);
            }
        \end{minted}
        \begin{minted}{c}
            void hanoi(int n, char a, char b, char c) {
                if (!n) {
                    return; // 递归出口
                }
                hanoi(n - 1, a, c, b);
                move(a, c);
                hanoi(n - 1, b, a, c);
            }
        \end{minted}
    }
\end{fragile}

\begin{frame}
    \frametitle{\insertsubsectionhead}
    \begin{alertblock}{注意}
        \begin{itemize}
            \item 递归可在抽象层面帮助快速理清思路
            \item 递归需借助栈结构存储额外信息，故效率较迭代/循环\alert{低}
            \item 在实现中，栈容量有限，故过多层递归极易\alert{爆栈}
        \end{itemize}
    \end{alertblock}
\end{frame}
% \begin{fragile}
    \frametitle{\insertsubsectionhead}
    \bicolumns[0.45]{
        \begin{exampleblock}{应用实例：迷宫寻路}
            \begin{itemize}
                \item 输入：迷宫
                \begin{itemize}
                    \item 由$\mathbb{R}^{m\times{}n}$上的矩阵表示
                    \item 矩阵元素可为\alert{墙壁}或\alert{地面}
                \end{itemize}
                \item 输出：自\alert{始}至\alert{终}的一条\alert{通路}
                \begin{itemize}
                    \item 通路由\alert{相邻}的地面序列组成
                    \item 相邻指\alert{八邻域}
                \end{itemize}
            \end{itemize}
        \end{exampleblock}
    }{
        \begin{figure}
            \centering
            \begin{bytefield}{2}
                \bitboxes{2}{{X}{\alert{始}}{X}{X}{X}{X}{X}{X}} \\
                \bitboxes{2}{{X}{\only<2>{$\checkmark$}}{X}{\only<2>{$\checkmark$}}{ }{ }{X}{X}} \\
                \bitboxes{2}{{X}{X}{\only<2>{$\checkmark$}}{X}{\only<2>{$\checkmark$}}{ }{ }{X}} \\
                \bitboxes{2}{{X}{X}{ }{X}{\only<2>{$\checkmark$}}{X}{X}{X}} \\
                \bitboxes{2}{{X}{X}{ }{X}{\only<2>{$\checkmark$}}{ }{ }{X}} \\
                \bitboxes{2}{{X}{ }{X}{X}{X}{\only<2>{$\checkmark$}}{X}{X}} \\
                \bitboxes{2}{{X}{X}{ }{X}{ }{\only<2>{$\checkmark$}}{X}{X}} \\
                \bitboxes{2}{{X}{X}{X}{X}{X}{X}{\alert{终}}{X}}
            \end{bytefield}
            \caption{迷宫示意图（X表示墙壁）}
            \label{fig:demo_maze}
        \end{figure}
    }
\end{fragile}

\begin{fragile}
    \frametitle{\insertsubsectionhead}
    \bicolumns[0.45]{
        \begin{exampleblock}{回溯法(Backtracking)思想}
            \begin{itemize}
                \item 一种试探纠错的搜索方法
            \end{itemize}
            \begin{enumerate}
                \item 自初始点出发
                \item 初始化搜索方向
                \item 按搜索方向向前搜索
                \item 判断搜索结果
                \begin{itemize}
                    \item 若到达终止点，则成功停止
                    \item 否则切换新搜索方向并转向步骤3
                \end{itemize}
                \item \alert{沿原路返回}上一结点并转向步骤2
                \item 失败停止
            \end{enumerate}
        \end{exampleblock}
    }{
        \begin{figure}
            \centering
            \begin{bytefield}{2}
                \bitboxes{2}{{X}{\alert{始}}{X}{X}{X}{X}{X}{X}} \\
                \bitboxes{2}{{X}{\only<2->{\alert{$\checkmark$}}}{X}{\only<12->{\alert{$\checkmark$}}}{ }{ }{X}{X}} \\
                \bitboxes{2}{{X}{X}{\only<3->{\alert{$\checkmark$}}}{X}{\only<13->{\alert{$\checkmark$}}}{ }{ }{X}} \\
                \bitboxes{2}{{X}{X}{\only<4->{\alert<4-10>{$\checkmark$}}}{X}{\only<14->{\alert{$\checkmark$}}}{X}{X}{X}} \\
                \bitboxes{2}{{X}{X}{\only<5->{\alert<5-9>{$\checkmark$}}}{X}{\only<15->{\alert{$\checkmark$}}}{ }{ }{X}} \\
                \bitboxes{2}{{X}{\only<6->{\alert<6-8>{$\checkmark$}}}{X}{X}{X}{\only<16->{\alert{$\checkmark$}}}{X}{X}} \\
                \bitboxes{2}{{X}{X}{\only<7->{\alert<7>{$\checkmark$}}}{X}{ }{\only<17->{\alert{$\checkmark$}}}{X}{X}} \\
                \bitboxes{2}{{X}{X}{X}{X}{X}{X}{\alert<18>{终}}{X}}
            \end{bytefield}
            \caption{回溯法一种可能的过程（X表示墙壁）}
            \label{fig:demo_backtracking}
        \end{figure}
    }
\end{fragile}

\begin{frame}
    \frametitle{\insertsubsectionhead}
    \begin{alertblock}{思考}
        \begin{itemize}
            \item 以何种数据结构表达迷宫？
            \item 如何表达寻路试探方向？
            \item 如何记录已试探结点以便原路返回？
            \item 如何避免重复试探相同结点？
        \end{itemize}
    \end{alertblock}
\end{frame}

\begin{fragile}
    \frametitle{\insertsubsectionhead}
    \bicolumns[0.45]{
        \begin{exampleblock}{表达迷宫的数据结构}
            \begin{itemize}
                \item 可采用二维整型数组
                \begin{itemize}
                    \item 0：地面
                    \item 1：障碍
                \end{itemize}
                \item 为实现方便可外加一圈障碍
            \end{itemize}
        \end{exampleblock}
        \begin{minted}{c}
            #define WIDTH 8
            #define HEIGHT 8

            static int map[HEIGHT + 2][WIDTH + 2];
        \end{minted}
    }{
        \begin{figure}
            \centering
            \begin{bytefield}{2}
                \bitboxes{2}{{1}{1}{1}{1}{1}{1}{1}{1}{1}{1}} \\
                \bitboxes{2}{{1}{1}{\alert{0}}{1}{1}{1}{1}{1}{1}{1}} \\
                \bitboxes{2}{{1}{1}{0}{1}{0}{0}{0}{1}{1}{1}} \\
                \bitboxes{2}{{1}{1}{1}{0}{1}{0}{0}{0}{1}{1}} \\
                \bitboxes{2}{{1}{1}{1}{0}{1}{0}{1}{1}{1}{1}} \\
                \bitboxes{2}{{1}{1}{1}{0}{1}{0}{0}{0}{1}{1}} \\
                \bitboxes{2}{{1}{1}{0}{1}{1}{1}{0}{1}{1}{1}} \\
                \bitboxes{2}{{1}{1}{1}{0}{1}{0}{0}{1}{1}{1}} \\
                \bitboxes{2}{{1}{1}{1}{1}{1}{1}{1}{\alert{0}}{1}{1}} \\
                \bitboxes{2}{{1}{1}{1}{1}{1}{1}{1}{1}{1}{1}}
            \end{bytefield}
            \caption{迷宫数据结构}
            \label{fig:data_structure_maze}
        \end{figure}
    }
\end{fragile}

\begin{fragile}
    \frametitle{\insertsubsectionhead}
    \bicolumns[0.55]{
        \begin{exampleblock}{寻路的试探方向}
            \begin{itemize}
                \item 可采用二维向量
            \end{itemize}
        \end{exampleblock}
        \begin{figure}
            \centering
            \begin{bytefield}{2}
                \bitboxes{8}{{$(x-1,y-1)$}{$(x+0,y-1)$}{$(x+1,y-1)$}} \\
                \bitboxes{8}{{$(x-1,y+0)$}{$(x+0,y+0)$}{$(x+1,y+0)$}} \\
                \bitboxes{8}{{$(x-1,y+1)$}{$(x+0,y+1)$}{$(x+1,y+1)$}}
            \end{bytefield}
            \caption{试探方向}
            \label{fig:searching_directions}
        \end{figure}
    }{
        \begin{minted}{c}
            #define N_DIRECTIONS 8

            typedef struct {
                int x, y;
            } Direction;

            static const Direction steps[N_DIRECTIONS] = {
                {1, 0}, {1, 1}, {0, 1}, {-1, 1},
                {-1, 0}, {-1, -1}, {0, -1}, {1, -1},
            };
        \end{minted}
    }
\end{fragile}

\begin{fragile}
    \frametitle{\insertsubsectionhead}
    \bicolumns[0.5]{
        \begin{exampleblock}{结点记录与原路返回}
            \begin{itemize}
                \item 可采用\alert{栈}记录已访问结点
                \item 结点信息须包含\alert{位置}与\alert{方向}
            \end{itemize}
        \end{exampleblock}
    }{
        \begin{minted}{c}
            typedef struct {
                int x; // 横坐标
                int y; // 纵坐标
                int d; // 方向数组下标
            } Cell; // 结点类型
        \end{minted}
    }
\end{fragile}

\begin{frame}
    \frametitle{\insertsectionhead：回到编辑器撤销}
    \begin{exampleblock}{用栈保存操作}
        从栈底到栈顶依次压入$A,B,C$；两次出栈得到$C,B$，
        分别撤销对应操作，栈中剩下$A$。
    \end{exampleblock}
    实际编辑器还需保存足以恢复状态的信息；空栈时不能再撤销。
\end{frame}

\begin{frame}
    \frametitle{\insertsectionhead：引例二}
    \begin{exampleblock}{打印任务排队}
        任务$A,B,C$依次到达一台打印机，打印机一次处理一个任务。
        如何让先到的任务先被处理？
    \end{exampleblock}
    \begin{block}{关键需求}
        新任务在末尾等待，处理时取最早的等待任务，需要先进先出。
    \end{block}
\end{frame}

\subsection{队列}

\begin{fragile}
    \frametitle{\insertsubsectionhead}
    \bicolumns[0.6]{
        \begin{block}{队列(Queue)}
            \begin{itemize}
                \item 另一种操作\alert{受限}的序列
                \item \alert{仅}允许在一端插入且在另一端删除元素
                \item 可插入的一端为\textbf{队尾(rear)}，另一端为\textbf{队首(front)}
                \item \textbf{先进先出(First In First Out, FIFO)}
            \end{itemize}
        \end{block}
        \uncover<2>{
            \begin{exampleblock}{日常生活中的队列}
                \begin{itemize}
                    \item 排队的人（不允许插队）
                    \item 羽毛球桶
                    \item $\dots$
                \end{itemize}
            \end{exampleblock}
        }
    }{
        \begin{figure}
            \centering
            \begin{bytefield}{6}
                \stackbox{3.5}{6}{$\vdots$\\\vspace*{1ex}$\downarrow$} \\
                \begin{leftwordgroup}{$n$}
                    \begin{rightwordgroup}{\mintinline[fontsize=\smaller]{c}{队尾}}
                        \stackbox{1.5}{6}{$a_{n}$}
                    \end{rightwordgroup} \\
                    \stackbox{2.5}{6}{$\vdots$\vspace{1.5ex}} \\
                    \stackbox{1.5}{6}{$a_{2}$} \\
                    \begin{rightwordgroup}{\mintinline[fontsize=\smaller]{c}{队首}}
                        \stackbox{1.5}{6}{$a_{1}$}
                    \end{rightwordgroup}
                \end{leftwordgroup} \\
                \stackbox{3.5}{6}{$\downarrow$\vspace*{1ex}\\$\vdots$}
            \end{bytefield}
            \caption{队列的示意图}
            \label{fig:demo_queue}
        \end{figure}
    }
\end{fragile}

\begin{fragile}
    \frametitle{\insertsectionhead}
    \begin{block}{队列的抽象数据类型}
        \begin{minted}[linenos=false,escapeinside=@@]{c}
            ADT Queue {
            数据:
                数据对象: @$\mathcal{D} = \{a_{k} | a_{k}\in\text{数据元素集合}, k\in\mathbb{Z}\cap[1,n]\}$@
                逻辑关系: @$\mathcal{R} = \{\langle{}a_{k-1},a_{k}\rangle | k\in\mathbb{Z}\cap[2,n]\}$@
            操作:
                create_queue()
                    构造并初始化一个空队列@$s$@
                is_empty_queue()
                    判断队列是否为空
                in_queue(s, x)
                    若队列@$s$@存在, 则在其队尾插入值为@$x$@的新元素
                out_queue(s)
                    若队列@$s$@存在且非空, 则记录并删除其队首元素
            }
        \end{minted}
    \end{block}
\end{fragile}

\begin{fragile}
    \frametitle{\insertsubsectionhead}
    \bicolumns[0.5]{
        \begin{block}{顺序队列（版本甲）}
            \begin{itemize}
                \item 可采用已有序列结构的封装
                \item 序列第一个元素为\alert{队首}
                \item 序列最后一个元素为\alert{队尾}
                \item 入队时可直接在末尾插入新元素
                \item 出队时可删除第一个元素
            \end{itemize}
        \end{block}
    }{
        \begin{minted}[highlightlines={}]{c}
            typedef struct {
                SequenceList *_;
            } SequenceQueue;
        \end{minted}
        \begin{minted}[highlightlines={}]{c}
            bool push_sequence_queue(
                    SequenceQueue *s, DataType d) {
                return s && insert_sequence_list(
                        s->_, s->_->last + 1, d);
            }
        \end{minted}
        \begin{minted}[highlightlines={}]{c}
            bool pop_sequence_queue(
                    SequenceQueue *s, DataType *d) {
                return s && remove_sequence_list(s->_, 0, d);
            }
        \end{minted}
    }
\end{fragile}

\begin{frame}
    \frametitle{\insertsubsectionhead}
    \begin{block}{时间复杂度}
        \begin{itemize}
            \item 入队操作：无需移动任何元素，$O(1)$
            \item 出队操作：需移动所有元素，\alert{$O(n)$，如何改进？}
        \end{itemize}
    \end{block}
    \pause
    \begin{alertblock}{思路}
        \begin{itemize}
            \item \alert{与其移动元素，不如移动队首，正如队尾一样}
        \end{itemize}
    \end{alertblock}
\end{frame}

\begin{fragile}
    \frametitle{\insertsubsectionhead}
    \bicolumns[0.5]{
        \begin{block}{顺序队列（版本乙）}
            \begin{itemize}
                \item 总体同版本甲
                \item 新增队首序号成员
                \item 出队时仅需增加队首序号
            \end{itemize}
        \end{block}
    }{
        \begin{minted}[highlightlines={}]{c}
            typedef struct {
                SequenceList *_;
                int front;
            } SequenceQueue;
        \end{minted}
        \begin{minted}[highlightlines={}]{c}
            bool pop_sequence_queue(
                    SequenceQueue *s, DataType *d) {
                if (!s) {
                    return false;
                }
                if (d) {
                    *d = s->_->data[front];
                }
                return ++front;
            }
        \end{minted}
    }
\end{fragile}

\begin{fragile}
    \frametitle{\insertsubsectionhead}
    \begin{block}{假上溢现象}
        \begin{itemize}
            \item 队首前尚有可用空间，但队尾已至序列尽头
        \end{itemize}
    \end{block}
    \begin{figure}
        \centering
        \begin{bytefield}{2}
            \bitboxes[]{3}{{}{}{\mintinline[fontsize=\smaller]{c}{front}}} \\
            \bitboxes[]{3}{{}{}{\tikzmark{ska}}} \\
            \bitbox[]{3}{$s:$}
            \bitboxes{3}{{\only<5->{$54$}}{}{$34$}{\only<2->{$23$}}{\only<3->{$67$}}{\only<4->{$43$}}} \\
            \bitboxes[]{3}{{}{\only<5>{\tikzmark{skb}}}{}{\only<1>{\tikzmark{skb}}}{\only<2>{\tikzmark{skb}}}{\only<3>{\tikzmark{skb}}}{\only<4>{\tikzmark{skb}}}} \\
            \bitboxes[]{3}{{}{\only<5>{\mintinline[fontsize=\smaller]{c}{rear}}}{}{\only<1>{\mintinline[fontsize=\smaller]{c}{rear}}}{\only<2>{\mintinline[fontsize=\smaller]{c}{rear}}}{\only<3>{\mintinline[fontsize=\smaller]{c}{rear}}}{\only<4>{\mintinline[fontsize=\smaller]{c}{rear}}}}
        \end{bytefield}
        \begin{tikzpicture}[ >=stealth, thick, black!50, overlay, remember picture, %
                list item/.style={draw=gray, rounded corners, thick} %
            ]
            \draw[->] ($({pic cs:{ska}})+(0,0.3)$) to [out=-90,in=90] ($({pic
                cs:{ska}})+(0,-0.3)$);
            \draw[<-] ($({pic cs:{skb}})+(0,0.25)$) to [out=-90,in=90] ($({pic
                cs:{skb}})+(0,-0.35)$);
        \end{tikzpicture}
        \caption{队列的假上溢现象及其解决思路}
        \label{fig:queue_fake_overflow}
    \end{figure}
    \pause
    \uncover<5>{
        \begin{alertblock}{思路}
            \begin{itemize}
                \item \alert{采用\textbf{循环队列}，假上溢时可将队尾移至内置数组
                          首元素，以充分利用空间}
            \end{itemize}
        \end{alertblock}
    }
\end{fragile}

\begin{fragile}
    \frametitle{\insertsubsectionhead}
    \begin{block}{新问题：如何判断队列已空或已满？}
        \vspace*{2ex}
        \bicolumns[0.5]{
            \begin{figure}
                \centering
                \begin{bytefield}{2}
                    \bitboxes[]{3}{{}{}{\mintinline[fontsize=\smaller]{c}{front}}} \\
                    \bitboxes[]{3}{{}{}{\tikzmark{skc}}} \\
                    \bitbox[]{3}{$s:$}
                    \bitboxes{3}{{}{}{}{}{}} \\
                    \bitboxes[]{3}{{}{}{\tikzmark{skd}}} \\
                    \bitboxes[]{3}{{}{}{\mintinline[fontsize=\smaller]{c}{rear}}}
                \end{bytefield}
                \begin{tikzpicture}[ >=stealth, thick, black!50, overlay, remember picture, %
                        list item/.style={draw=gray, rounded corners, thick} %
                    ]
                    \draw[->] ($({pic cs:{skc}})+(0,0.3)$) to [out=-90,in=90] ($({pic
                        cs:{skc}})+(0,-0.3)$);
                    \draw[<-] ($({pic cs:{skd}})+(0,0.25)$) to [out=-90,in=90] ($({pic
                        cs:{skd}})+(0,-0.35)$);
                \end{tikzpicture}
                \caption{队列已空}
                \label{fig:queue_empty}
            \end{figure}
        }{
            \begin{figure}
                \centering
                \begin{bytefield}{2}
                    \bitboxes[]{3}{{}{}{\mintinline[fontsize=\smaller]{c}{front}}} \\
                    \bitboxes[]{3}{{}{}{\tikzmark{ske}}} \\
                    \bitbox[]{3}{$s:$}
                    \bitboxes{3}{{$54$}{$34$}{$23$}{$67$}{$43$}} \\
                    \bitboxes[]{3}{{}{}{\tikzmark{skf}}} \\
                    \bitboxes[]{3}{{}{}{\mintinline[fontsize=\smaller]{c}{rear}}}
                \end{bytefield}
                \begin{tikzpicture}[ >=stealth, thick, black!50, overlay, remember picture, %
                        list item/.style={draw=gray, rounded corners, thick} %
                    ]
                    \draw[->] ($({pic cs:{ske}})+(0,0.3)$) to [out=-90,in=90] ($({pic
                        cs:{ske}})+(0,-0.3)$);
                    \draw[<-] ($({pic cs:{skf}})+(0,0.25)$) to [out=-90,in=90] ($({pic
                        cs:{skf}})+(0,-0.35)$);
                \end{tikzpicture}
                \caption{队列已满}
                \label{fig:queue_full}
            \end{figure}
        }
    \end{block}
\end{fragile}

\begin{frame}
    \frametitle{\insertsubsectionhead}
    \begin{block}{思路甲：设立标志位\mintinline[fontsize=\smaller]{c}{flag}辅助判断}
        \begin{itemize}
            \item 除检查队首与队尾是否重合外，尚需判断标志位（可约定：0为队空，1
                  为队满）
        \end{itemize}
    \end{block}
    \pause
    \begin{block}{思路乙：在队首处空出一个元素}
        \begin{itemize}
            \item 队空时队首队尾重合，队满时二者不重合
        \end{itemize}
    \end{block}
    \pause
    \begin{block}{思路丙：新增\mintinline[fontsize=\smaller]{c}{size}成员记录队列长度}
        \begin{itemize}
            \item 队空时长度为0，队满时长度为容量
        \end{itemize}
    \end{block}
\end{frame}

\begin{fragile}
    \frametitle{\insertsubsectionhead}
    \bicolumns[0.5]{
        \begin{block}{顺序队列（版本丙）}
            \begin{itemize}
                \item 总体同版本乙
                \item 新增队列长度成员
                \item 因出入较大，故已不借用序列结构
            \end{itemize}
        \end{block}
    }{
        \begin{minted}[highlightlines={}]{c}
            typedef struct {
                DataType data[CAPACITY];
                int front; // 队首
                int rear; // 队尾
                int size; // 队列长度
            } SequenceQueue;
        \end{minted}
    }
\end{fragile}

\begin{fragile}
    \frametitle{\insertsubsectionhead}
    \bicolumns[0.5]{
        \begin{block}{顺序队列的初始化}
            \begin{itemize}
                \item 为队列分配空间
                \item 将队首与队尾均置于末尾
                \item 将队列长度置0
            \end{itemize}
        \end{block}
    }{
        \begin{minted}[highlightlines={}]{c}
            SequenceQueue *create_sequence_queue() {
                SequenceQueue *s = malloc(sizeof(SequenceQueue));
                if (s) {
                    s->front = s->rear = CAPACITY - 1;
                    s->size = 0;
                }
                return s;
            }
        \end{minted}
    }
\end{fragile}

\begin{fragile}
    \frametitle{\insertsubsectionhead}
    \bicolumns[0.5]{
        \begin{block}{入队}
            \begin{itemize}
                \item 确保队列未满
                \item 队尾后移
                      \begin{itemize}
                          \item 若已至绝境则从头再来
                      \end{itemize}
                \item 加入新元素
                \item 更新队列长度
            \end{itemize}
        \end{block}
    }{
        \begin{minted}[highlightlines={}]{c}
            bool push_sequence_queue(
                    SequenceQueue *s, DataType d) {
                if (full_sequence_queue(s)) {
                    return false;
                }
                s->rear = (s->rear + 1) % CAPACITY;
                s->data[s->rear] = d;
                return ++s->size;
            }
        \end{minted}
    }
\end{fragile}

\begin{fragile}
    \frametitle{\insertsubsectionhead}
    \bicolumns[0.5]{
        \begin{block}{出队}
            \begin{itemize}
                \item 确保队列非空
                \item 队首后移\footnotemark
                      \begin{itemize}
                          \item 若已至绝境则从头再来
                      \end{itemize}
                \item 记录已删除元素
                \item 更新队列长度
            \end{itemize}
        \end{block}
    }{
        \begin{minted}[highlightlines={}]{c}
            bool pop_sequence_queue(
                    SequenceQueue *s, DataType *d) {
                if (empty_sequence_queue(s)) {
                    return false;
                }
                s->front = (s->front + 1) % CAPACITY;
                if (d) {
                    *d = s->data[s->front];
                }
                --s->size;
                return true;
            }
        \end{minted}
    }

    \footnotetext{注：队首一般指向首元素\alert{前一个}位置}
\end{fragile}

\begin{fragile}
    \frametitle{\insertsubsectionhead}
    \bicolumns[0.5]{
        \begin{block}{链式队列}
            \begin{itemize}
                \item 可采用已有链表结构的封装
                \item 新增队首与队尾指针
            \end{itemize}
        \end{block}
    }{
        \begin{minted}[highlightlines={}]{c}
            typedef struct {
                LinkedList *_;
                LinkedNode *front; // 队首
                LinkedNode *rear; // 队尾
            } LinkedQueue;
        \end{minted}
    }
\end{fragile}

\begin{fragile}
    \frametitle{\insertsubsectionhead}
    \begin{exampleblock}{链式队列}
        \vspace*{2ex}
        \bicolumns[0.3]{
            \begin{figure}
                \centering
                \begin{bytefield}{2}
                    \bitboxes[]{3}{{}{\mintinline[fontsize=\smaller]{c}{front}}} \\
                    \bitboxes[]{3}{{}{\tikzmark{smf}}} \\
                    \bitbox[]{3}{$s:$}
                    \bitbox{3}{}\bitbox{2}{\tikzmark{sma}} \\
                    \bitbox[]{3}{}
                    \bitbox[]{3}{\tikzmark{smg}} \\
                    \bitbox[]{3}{}
                    \bitbox[]{3}{\mintinline[fontsize=\smaller]{c}{rear}}
                \end{bytefield}
                \begin{tikzpicture}[ >=stealth, thick, black!50, overlay, remember picture, %
                        list item/.style={draw=gray, rounded corners, thick} %
                    ]
                    \draw ($({pic cs:{sma}})+(-0.22,-0.3)$) -- ($({pic cs:{sma}})+(0.23,0.23)$);
                    \draw[->] ($({pic cs:{smf}})+(0,0.3)$) to [out=-90,in=90] ($({pic
                        cs:{smf}})+(0,-0.3)$);
                    \draw[<-] ($({pic cs:{smg}})+(0,0.25)$) to [out=-90,in=90] ($({pic
                        cs:{smg}})+(0,-0.35)$);
                \end{tikzpicture}
                \caption{空链式队列}
                \label{fig:demo_empty_linked_queue}
            \end{figure}
        }{
            \begin{figure}
                \centering
                \begin{bytefield}{2}
                    \bitboxes[]{3}{{}{\mintinline[fontsize=\smaller]{c}{front}}} \\
                    \bitboxes[]{3}{{}{\tikzmark{slf}}} \\
                    \bitbox[]{3}{$s:$}
                    \bitbox{3}{}\bitbox{2}{\tikzmark{sla}}\bitbox[]{2}{\tikzmark{tla}}
                    \bitbox{3}{$34$}\bitbox{2}{\tikzmark{slb}}\bitbox[]{2}{\tikzmark{tlb}}
                    \bitbox{3}{$23$}\bitbox{2}{\tikzmark{slc}}\bitbox[]{2}{\tikzmark{tlc}}
                    \bitbox{3}{$43$}\bitbox{2}{\tikzmark{sld}} \\
                    \bitbox[]{24}{}
                    \bitbox[]{3}{\tikzmark{slg}} \\
                    \bitbox[]{24}{}
                    \bitbox[]{3}{\mintinline[fontsize=\smaller]{c}{rear}}
                \end{bytefield}
                \begin{tikzpicture}[ >=stealth, thick, black!50, overlay, remember picture, %
                        list item/.style={draw=gray, rounded corners, thick} %
                    ]
                    \draw[->] ($({pic cs:{sla}})+(0,-0.03)$) to [out=0,in=180] ($({pic
                        cs:{tla}})+(0.26,-0.03)$);
                    \fill ($({pic cs:{sla}})+(0,-0.03)$) circle [radius=0.05];
                    \draw[->] ($({pic cs:{slb}})+(0,-0.03)$) to [out=0,in=180] ($({pic
                        cs:{tlb}})+(0.26,-0.03)$);
                    \fill ($({pic cs:{slb}})+(0,-0.03)$) circle [radius=0.05];
                    \draw[->] ($({pic cs:{slc}})+(0,-0.03)$) to [out=0,in=180] ($({pic
                        cs:{tlc}})+(0.26,-0.03)$);
                    \fill ($({pic cs:{slc}})+(0,-0.03)$) circle [radius=0.05];
                    \draw ($({pic cs:{sld}})+(-0.22,-0.3)$) -- ($({pic cs:{sld}})+(0.23,0.23)$);
                    \draw[->] ($({pic cs:{slf}})+(0,0.3)$) to [out=-90,in=90] ($({pic
                        cs:{slf}})+(0,-0.3)$);
                    \draw[<-] ($({pic cs:{slg}})+(0,0.25)$) to [out=-90,in=90] ($({pic
                        cs:{slg}})+(0,-0.35)$);
                \end{tikzpicture}
                \caption{非空链式队列}
                \label{fig:demo_non_empty_linked_queue}
            \end{figure}
        }
    \end{exampleblock}
\end{fragile}

\begin{fragile}
    \frametitle{\insertsubsectionhead}
    \bicolumns[0.5]{
        \begin{block}{链式队列的初始化}
            \begin{itemize}
                \item 调用链表初始化方法
                \item 将队首与队尾均指向首结点
            \end{itemize}
        \end{block}
    }{
        \begin{minted}[highlightlines={}]{c}
            LinkedQueue *create_linked_queue() {
                LinkedQueue *s = malloc(sizeof(LinkedQueue));
                if (s) {
                    s->_ = create_linked_list();
                    s->front = s->rear = s->_->head;
                }
                return s;
            }
        \end{minted}
    }
\end{fragile}

\begin{fragile}
    \frametitle{\insertsubsectionhead}
    \bicolumns[0.5]{
        \begin{block}{链式队列的判空}
            \begin{itemize}
                \item 队首与队尾是否重合
            \end{itemize}
        \end{block}
    }{
        \begin{minted}[highlightlines={}]{c}
            bool empty_linked_queue(LinkedQueue *s) {
                return !s || s->front == s->rear;
            }
        \end{minted}
    }
\end{fragile}

\begin{fragile}
    \frametitle{\insertsubsectionhead}
    \bicolumns[0.5]{
        \begin{block}{入队}
            \begin{itemize}
                \item 创建新结点
                \item 在链表末端插入新结点\footnotemark
                \item 队尾指针后移
            \end{itemize}
        \end{block}
    }{
        \begin{minted}[highlightlines={}]{c}
            bool push_linked_queue(
                    LinkedQueue *s, DataType d) {
                if (!s) {
                    return false;
                }
                LinkedNode *n = create_linked_node(d);
                return n && (s->rear =
                    attach_after_node(s->rear, n));
            }
        \end{minted}
    }

    \footnotetext{注：此时已有队尾指针记录队尾结点，故无需每次查找之}
\end{fragile}

\begin{fragile}
    \frametitle{\insertsubsectionhead}
    \bicolumns[0.5]{
        \begin{block}{出队}
            \begin{itemize}
                \item 记录链表第一结点的值
                \item 删除该首结点
                \item 当队列空时队尾指针前移
            \end{itemize}
        \end{block}
    }{
        \begin{minted}[highlightlines={}]{c}
            bool pop_linked_queue(
                    LinkedQueue *s, DataType *d) {
                if (empty_linked_queue(s)) {
                    return false;
                }
                LinkedNode *p = detach_after_node(s->front);
                if (d) {
                    *d = p->data;
                }
                free(p);
                if (empty_linked_queue(s)) {
                    s->rear = s->front;
                }
                return true;
            }
        \end{minted}
    }
\end{fragile}


\begin{frame}
    \frametitle{\insertsectionhead：回到打印任务}
    \begin{exampleblock}{用队列安排任务}
        $A,B,C$依次入队，随后出队$A$开始打印；此时$D$到达并入队。
        余下任务从队首到队尾为$B,C,D$，继续依次处理。
    \end{exampleblock}
    队列为空时没有等待任务；任务取消、优先级等需求需要额外规则。
\end{frame}

\begin{frame}
    \frametitle{\insertsectionhead}
        \begin{table}
            \small
            \caption{一点建议}
            \label{tab:stack_queue_suggestions}
            \begin{tabular}{rcc}
                \toprule
                & \textbf{栈} & \textbf{队列} \\
                \midrule
                \textbf{应用} & 逆序记录 & 顺序缓冲 \\
                \textbf{实现} & 动态数组 & 双向链表 \\
                \bottomrule
            \end{tabular}
        \end{table}
\end{frame}